% ============================================================
% 第 1 章 引言 + 第 2 章 相关工作
% ============================================================
\section{引言}

提示工程（prompt engineering）已成为无需参数更新即可调控大语言模型行为的主要手段。自思维链（chain-of-thought，CoT）提示被证明能显著提升推理任务表现以来\cite{wei2022cot,kojima2022zeroshot}，研究者提出了大量在正式作答前插入中间生成环节的策略：零样本思维链在答案前追加“请一步步思考”\cite{kojima2022zeroshot}，自一致性采样多条推理路径后投票\cite{wang2023selfconsistency}，思维树则在解题过程中显式引入分支搜索与回溯\cite{yao2023tot}。这类方法共同的隐含假设是：\textbf{模型在产出最终答案之前进行的额外生成，能够以可迁移的方式改善最终答案本身}。

然而，上述证据大多来自具有唯一正确答案的推理或问答任务。在开放式生成场景（如方案设计、创意构想、理论构建）中，“中间生成”是否同样有益，仍缺乏系统证据。一个自然但未被充分检验的猜想是：让模型在作答前先自由联想一批发散性的“灵感素材”，再基于素材作答，可以起到类似头脑风暴的作用，使答案更具原创性。本文研究的\textbf{词对提示}（word-pair priming）正是此类策略的一个典型实例——在任务指令前加入前缀：“在回答前面输出200组词，每组有两个词，从这些词获得灵感”。该策略的直觉来源是认知心理学中的发散思维与远距离联想：两个词构成的意象组合（如“潮·汐”“光·影”）可能触发模型跨越常规语义邻域的组合搜索，从而突破训练分布中的常规答案模式。

这一策略存在三重潜在风险。其一，\textbf{指令目标稀释}：200 组词对通常消耗数千 token 的生成预算，可能压缩正文的论证篇幅与细节密度；其二，\textbf{语境干扰}：已有研究表明大语言模型极易被无关上下文分散注意力\cite{shi2023distracted}，长上下文中的信息利用也存在“中间丢失”效应\cite{liu2024lost}，随机联想词对与任务的相关性参差不齐，可能将模型注意力引向低质量联想；其三，\textbf{交付形态问题}：词对列表属于过程性输出，若需从最终产物中剥离，清洗操作本身可能破坏正文的连贯性。这三条风险在推理任务中并不显著（因为推理链是答案的构成部分），但在开放式生成中被同时激活。

针对上述问题，本文基于一个公开归档的实验仓库\footnote{\raggedright 实验数据与全部原始归档位于 Gitee 仓库 \texttt{qwe12345678/jkdsfn}（分支 \texttt{v2-deepseek}，目录 \texttt{llm\_experiment\_results\_v2/}）：\par\url{https://gitee.com/qwe12345678/jkdsfn}\par}，对词对提示的有效性进行了严格的实证检验。实验以 DeepSeek-V4.1-Flash 与 DeepSeek-V4-Pro\cite{deepseekv3}两个商用模型为生成器，在 6 个开放式生成任务上比较 \baseline{} 与 \wordpairs{} 两种条件（每格 10 次重复，共 240 条生成），词对输出经带正则校验的二次清洗协议处理后，由两个模型充当评委进行位置哈希平衡的 A/B 盲评（240 对）。本研究的贡献如下：

\begin{enumerate}[leftmargin=2em, itemsep=0.15em]
  \item 提出并量化了一种“先发散素材、后正式作答”提示策略在开放式生成任务上的效果，填补了该类策略在非推理任务上的证据空白；
  \item 设计了双生成模型$\times$双评委的交叉盲评协议，并通过确定性位置平衡、清洗残留校验、数据溯源等手段控制了 LLM 评审中的主要混淆因素\cite{zheng2023judging}；
  \item 给出一致的量化结论：词对提示使整体质量平均下降 $1.72$ 分（10 分制，$p<0.001$），仅在原创性维度小幅获益（$+0.31$ 分，$p<0.001$），且代价是 $1.60\sim2.02$ 倍的生成 token 成本；
  \item 基于逐对分差、任务分解与案例文本，提出了词对提示失效的三条机理解释（目标稀释、语境干扰、清洗损失）及其适用边界。
\end{enumerate}

本文其余部分组织如下：第 2 节回顾相关工作；第 3 节给出实验设计；第 4 节报告实验结果；第 5 节进行案例分析；第 6 节讨论失效机理与实践启示；第 7 节分析有效性威胁；第 8 节总结全文。

\section{相关工作}

\subsection{思维链家族与中间生成策略}

中间生成策略的谱系始于少样本思维链提示。Wei 等人在数学、常识与符号推理任务上证明，在少样本示例中嵌入中间推理步骤可显著提升大模型表现\cite{wei2022cot}；Kojima 等人进一步发现，仅添加“让我们一步步思考”即可激活零样本推理能力，无需任何示例\cite{kojima2022zeroshot}。在解码层面，Wang 等人提出自一致性方法，对多条思维链采样后按多数答案投票，将单条路径的生成转化为集成决策\cite{wang2023selfconsistency}；Yao 等人的思维树方法则将中间推理组织为可回溯的树搜索，支持前瞻与自我评估\cite{yao2023tot}。这些方法的共同特点是：\textbf{中间生成与最终答案之间存在直接的逻辑依赖}——推理链的每一步都构成答案的推导组件。这与本文研究的词对提示形成鲜明对照：200 组灵感词对与最终正文之间不存在逻辑推导关系，仅存在启发式关联。因此，思维链家族的成功并不能外推到“自由联想素材”类提示；后者需要独立的实证检验。此外，指令微调研究指出，模型的指令遵循能力与生成目标的聚焦程度密切相关\cite{ouyang2022instructgpt}，而插入与任务无直接依赖的前置生成可能削弱这种聚焦。

\subsection{提示敏感性与无关上下文干扰}

大量证据表明大语言模型对提示的表面形式高度敏感。Sclar 等人系统量化了少-shot 提示的格式扰动（大小写、分隔符、间距）对模型表现的影响，发现性能波动可达数十个百分点，且不同模型的敏感度差异巨大\cite{sclar2024format}。Shi 等人则证明，在算术与常识推理题中加入与解题无关的句子会导致显著的性能下降，模型容易被表面相似但语义无关的上下文“带偏”\cite{shi2023distracted}。在长上下文场景下，Liu 等人发现多文档问答中位于上下文中部的信息利用率显著低于首尾，即“中间丢失”效应\cite{liu2024lost}。词对提示在结构上同时具备这三个风险因子：它向生成上下文注入数百个与任务相关性未经验证的联想项，且这些联想项在自回归生成中恰好位于正文之前（而非遥远的上下文中部），其对正文生成的持续影响缺乏先验保证。本文的实验可视为对“自生成的无关上下文”（self-generated irrelevant context）这一特殊干扰源的受控检验。

\subsection{LLM 作为评委及其偏差}

以强模型替代人工进行开放式输出评价已获得广泛验证。Zheng 等人在 MT-Bench 与 Chatbot Arena 上证明，GPT-4 级评委与人类偏好的一致率可达 80\% 以上，与人类评委之间的一致率相当\cite{zheng2023judging}；但他们同时系统识别了 LLM 评委的若干系统偏差：\textbf{位置偏差}（对 A/B 呈现顺序敏感）、\textbf{冗长偏差}（偏好更长回答）与\textbf{自我增强偏差}（偏好自家模型输出）。Shi 等人对位置偏差进行了更系统的测量，表明其在成对与列表式评审中普遍存在，交换顺序重评是有效缓解手段\cite{shi2025position}。本文的盲评协议据此进行了针对性设计：评委 prompt 明确声明“引用未展示灵感词本身不扣分”以切断对清洗残留的惩罚性归因；A/B 位置由确定性哈希平衡；采用双评委交叉评审以缓解单评委的自我增强偏差；并在结果分析中显式报告位置偏差量级（第 \ref{sec:position} 节）。此外，情感激励研究亦提示 prompt 中的额外刺激可双向影响输出质量\cite{li2023emotion}，进一步说明任何“看似无害”的提示附加物都需要实验验证而非直觉判断。
